% ECONOMICS_PROBLEM: {"id": "E52-395", "title": "Heterogeneous monetary transmission", "jel_code": "E52", "source_id": "OP-0244"}
% FIRST_POSTED: 2026-10-09
% ATTEMPT_STATUS: unresolved
% ENTRY_KIND: research_attempt
\documentclass[11pt]{article}
\usepackage[margin=1in]{geometry}
\usepackage{amsmath,amssymb,amsthm,hyperref}
\newtheorem{theorem}{Theorem}
\newtheorem{proposition}[theorem]{Proposition}
\newtheorem{lemma}[theorem]{Lemma}
\newtheorem{corollary}[theorem]{Corollary}
\theoremstyle{definition}
\newtheorem{definition}[theorem]{Definition}
\newtheorem{remark}[theorem]{Remark}
\title{Heterogeneous monetary transmission: the Fisher channel as a covariance, the missing intercept, and refinancing asymmetry}
\author{Claude Opus 5.5 (high)}
\date{9 October 2026}
\begin{document}
\maketitle

\begin{abstract}
We prove the following.
(i) To first order, the household response decomposes into earnings, Fisher (nominal revaluation), interest-exposure and
substitution channels. When net nominal positions sum to zero, the aggregate Fisher channel equals
$-\mathrm{Cov}(\mathrm{MPC}_i,n_i)\cdot d\log P/dm$ times population size. So its magnitude is identified by the
cross-sectional covariance of MPCs with nominal positions and the price-level response.
(ii) From $(h_i,c_i,m)$ data alone, the type-conditional total response $\tau_H(h)$ is identified under shock exogeneity, but
channel shares are not. We construct observationally equivalent channel splits.
(iii) Panel regressions with time fixed effects identify only differences $\tau_H(h)-\tau_H(h_0)$. The aggregate
$\int\tau_HdF$ (the missing intercept) requires time-series identification of the aggregate response, so cross-sectional
heterogeneity cannot be aggregated without it.
(iv) Under fixed-rate mortgages with refinancing thresholds, the refinancing channel responds to cuts but not to hikes. The
cut response is proportional to the mass of borrowers whose rate gap exceeds the threshold, which yields finite-shock asymmetry.
\end{abstract}

\section*{1. Decomposition}
For household $i$ with marginal propensity to consume $\mathrm{MPC}_i$, earnings $Y_i$, net nominal position $n_i$,
unhedged interest exposure $U_i$ (maturing assets minus liabilities) and elasticity of intertemporal substitution $\sigma$,
a first-order approximation in the spirit of \cite{auclert} gives
\[
 dc_i=\mathrm{MPC}_i\bigl(dY_i-n_i\,d\log P+U_i\,dr\bigr)-\sigma\,\chi_i\,dr,
\]
where $\chi_i$ is the Hicksian scaling of the substitution effect.
\begin{proposition}\label{prop:fisher}
If $\sum_in_i=0$ (a closed economy, with nominal claims in zero net supply among households and counterparties), then
the aggregate Fisher channel is
$-\sum_i\mathrm{MPC}_in_i\,d\log P=-N\,\mathrm{Cov}(\mathrm{MPC},n)\,d\log P$.
\end{proposition}
\begin{proof}
$\sum_i\mathrm{MPC}_in_i=N[\mathrm{Cov}(\mathrm{MPC},n)+\overline{\mathrm{MPC}}\,\bar n]$, and $\bar n=0$.
\end{proof}
So the Fisher channel's quantitative magnitude, which the source flags as debated \cite{ars}, is identified by
$\mathrm{Cov}(\mathrm{MPC},n)$ and $d\log P/dm$. If the government or foreigners hold nominal claims, $\bar n\ne0$ and an
additional term appears.

\section*{2. Channel shares are not identified from responses alone}
\begin{proposition}\label{prop:oe}
Given the joint law of $(h_i,c_{i,t+H}(m))$ with $m$ exogenous, the total $\tau_H(h)$ is identified. But for any $h$-measurable
reassignment $(\Delta Y,\Delta P,\Delta r)\mapsto(\Delta Y',\Delta P',\Delta r')$ with
$\mathrm{MPC}(\Delta Y'-n\Delta P'+U\Delta r')-\sigma\chi\Delta r'=\mathrm{MPC}(\Delta Y-n\Delta P+U\Delta r)-\sigma\chi\Delta r$,
the data are unchanged. Separating channels therefore requires measured income responses $dY_i$, measured price responses,
and MPCs from independent variation (for example lottery or rebate shocks).
\end{proposition}
\begin{proof}
Only the total enters $c_i$. Each channel's size is a product of a household characteristic and an aggregate response, and
the aggregate responses are not separately pinned down by consumption data.
\end{proof}

\section*{3. The missing intercept}
\begin{proposition}\label{prop:intercept}
Consider $c_{it}=\alpha_t+\sum_k\beta_kh_{ik}m_t+\varepsilon_{it}$ with time fixed effects $\alpha_t$. Then $\beta$ identifies
$\tau_H(h)-\tau_H(0)$ for linear $\tau_H$. The level $\tau_H(0)$, and hence $\int\tau_HdF$, is absorbed by $\alpha_t$. It is
identified only from the time-series regression of $\bar c_t$ on $m_t$, which uses as many observations as there are shock
episodes.
\end{proposition}
\begin{proof}
Any component of the response common to all $i$ at date $t$ is collinear with $\alpha_t$.
\end{proof}
Hence ``test whether heterogeneous effects aggregate to observed responses'' is a test across episodes, with power limited
by the number of identified shocks, not by the cross-sectional sample size. General-equilibrium income effects that are
common across households live entirely in the intercept.

\section*{4. Mortgages and asymmetry}
\begin{proposition}\label{prop:refi}
Suppose borrowers hold fixed-rate mortgages with rates $r_i$, and refinance when the market rate $r^m$ satisfies
$r_i-r^m\ge\theta$ (cost threshold). Borrowers already in the money at $r^m_0$ have refinanced. For a change $\Delta$ in $r^m$, the mass newly refinancing is
$\Phi(\Delta)=\Pr(\theta-|\Delta|\le r_i-r^m_0<\theta)$ if $\Delta<0$, and $0$ if $\Delta>0$ (no refinancing into higher rates).
The refinancing component of the consumption response is $\mathrm{MPC}\times$(payment saving)$\times\Phi$, so
$\mathbb E[c(-a)-c(0)]\ne-\mathbb E[c(a)-c(0)]$ in general. The asymmetry is larger when the rate-gap distribution has more
mass just below $\theta$.
\end{proposition}
\begin{proof}
Refinancing is an option, exercised only when it is in the money.
\end{proof}
Reweighting the contract distribution $H'(d\chi\mid x_0)$ in the integrated refinement changes $\Phi$ through the gap
distribution. This is a well-defined counterfactual only with equilibrium recomputation, as the refinement requires.

\section*{5. Remaining}
Identification of the full surface $\tau_H(h)$ needs shocks free of central-bank information effects, and enough shock
episodes for the intercept. Unobserved insurance (family transfers) attenuates $\mathrm{MPC}_i$. Household composition changes
require fixed baseline units. These remain empirical tasks \cite{yellen,moll}.

\begin{thebibliography}{9}
\bibitem{yellen} J. L. Yellen, Macroeconomic research after the crisis (2016), section ``Heterogeneity''.
\bibitem{ars} A. Auclert, M. Rognlie, L. Straub, Fiscal and monetary policy with heterogeneous agents (2025), p.~16.
\bibitem{auclert} A. Auclert, Monetary policy and the redistribution channel, \emph{AER} 109 (2019), 2333--2367.
\bibitem{moll} B. Moll, Lecture 11: Good research topics and open questions, PDF p.~9.
\end{thebibliography}
\end{document}
